\begin{frame}[fragile]{Main Idea of Promises}

\begin{itemize}
	\item A Promise is an \underbar{object} that (a) \underbar{represents a task} with two possible outcomes (success or failure) and (b)\underbar{holds callbacks} that fire when one outcome or the other has occurred.
	\item This is an example of jQuery code that uses callback and promises. 
\end{itemize}

\begin{columns}%
\begin{column}{0.4\textwidth}%
\begin{Code}{Callback}%
\begin{lstlisting}[firstnumber=1, label=glabels, xleftmargin=0pt, basicstyle=\scriptsize]% 

$.get('/mydata', { 
  success: onSuccess, 
  failure: onFailure, 
  always: onAlways
}); 
\end{lstlisting}%   
\end{Code}%
\end{column}%
\begin{column}{0.5\textwidth}%
\begin{Code}{Promises}%
\begin{lstlisting}[firstnumber=1, label=glabels, xleftmargin=0pt, basicstyle=\scriptsize]% 

var promise = 
  $.get('/mydata'); 
promise.done(onSuccess); 
promise.fail(onFailure); 
promise.always(onAlways);
\end{lstlisting}%   
\end{Code}%
\end{column}%
\end{columns}%
\end{frame}

\begin{frame}[fragile]{}

\begin{itemize}
	\item We can understand that Promises are nothing more than syntactic sugar to use callback functions easily.
	\item JavaScript even loop manages the callback queue and the thread pool in the same way both for callbacks and promises. 
\end{itemize}

\begin{center}
\includegraphics[width=0.5\textwidth]{javascript-event-loop.png}
\end{center} 
\end{frame}

\begin{frame}[fragile]{What is Promise in Node.js?}

\begin{itemize}
	\item For the Promises of JavaScript with web browsers, we often use jQuery's promises. 
	\item We focus on discussing the Promise of JavaScript/Node.js. 
	\item This code shows a template of a Node.js function that uses the Promise object.
	\item Notice that the Promise uses a callback function with two parameters, resolve and reject (lines 2 - 7).
\end{itemize}

\begin{Code}{}%
\begin{lstlisting}[firstnumber=1, label=glabels, xleftmargin=0pt, basicstyle=\tiny]% 

function A() {
 return new Promise((resolve, reject) => {
    ...
    if (...) resolve(...)
    else     reject(...)
 }  
)}
\end{lstlisting}%   
\end{Code}%

\end{frame}

\begin{frame}[fragile]{}
% https://blog.logrocket.com/guide-promises-node-js/
\begin{itemize}
	\item This is an example Node.js function getSumNum that promises to return the sum of two input values, a and b (it may take a while to get the result). 
	\item When the computation is finished, the result will be returned by invoking the resolve() function (when successful) or reject() function (when failed). 
\end{itemize}

\begin{Code}{}%
\begin{lstlisting}[firstnumber=1, label=glabels, xleftmargin=0pt, basicstyle=\scriptsize]% 

function getSumNum(a, b) {
  const customPromise = new Promise((resolve, reject) => {
    const sum = a + b;
    if(sum <= 5){
      resolve(sum)
    } else {
      reject(new Error('Oops!.. Number must be less than 5'))
    }
  })
  return customPromise
}
\end{lstlisting}%   
\end{Code}%
\end{frame}

\begin{frame}[fragile]{}

\begin{itemize}
	\item The promise object that is returned (from getSumNum(1,3)) has then() method that connects to the resolve() and catch() method that connects to the reject(). 
	\item In other words, when the promise object finishes its computation, it invokes the resolve (then) method to return the sum, and when there is an error, it invokes the reject (catch) method. 
\end{itemize}

\begin{Code}{}%
\begin{lstlisting}[firstnumber=1, label=glabels, xleftmargin=0pt, basicstyle=\scriptsize]% 

getSumNum(1, 3)
.then(data => {
  console.log("initial data: " + data)
  value = data + 1 
  console.log("received data: " + value)
})
.catch(err => {
  console.log(err)
})
\end{lstlisting}%   
\end{Code}%

\end{frame}